\subsubsection{Smer reťazca: 5'$\rightarrow$3' a 3'$\rightarrow$5'}
DNA reťazec má chemickú polaritu --- jeden koniec reťazca má voľnú fosfátovú skupinu na 5' uhlíku cukru (5' koniec) a druhý koniec má voľnú hydroxylovú skupinu na 3' uhlíku (3' koniec).

\begin{itemize}
    \item \textbf{Smer 5'$\rightarrow$3'} --- smer, v ktorom DNA polymeráza syntetizuje nový reťazec. Polymeráza číta templát v smere 3'$\rightarrow$5' a zapisuje nový reťazec v smere 5'$\rightarrow$3'.
    \item \textbf{Smer 3'$\rightarrow$5'} --- smer čítania templátového vlákna počas syntézy. Nie je možné syntetizovať reťazec v tomto smere priamo.
\end{itemize}

Táto polarita má priamy dopad na nanopórové sekvenovanie --- tá istá DNA molekula môže vstúpiť do nanopóru z ľubovoľného konca, čím vznikajú čítania v oboch smeroch (forward aj reverse). Pri spracovaní dát je preto potrebné uvažovať s oboma orientáciami.
